% Generated draft: requires mathematical and source review.
\documentclass[11pt]{article}
\usepackage[margin=25mm]{geometry}
\usepackage{fontspec}
\setmainfont{Latin Modern Roman}
\usepackage{amsmath,amssymb,mathtools,unicode-math}
\setmathfont{Latin Modern Math}
\usepackage{enumitem,needspace,xurl,hyperref}
\setlength{\parindent}{0pt}
\setlength{\parskip}{5pt}
\setlength{\emergencystretch}{3em}
\newcommand{\N}{\mathbb{N}}
\newcommand{\Z}{\mathbb{Z}}
\newcommand{\Q}{\mathbb{Q}}
\newcommand{\R}{\mathbb{R}}
\newcommand{\C}{\mathbb{C}}
\newcommand{\F}{\mathbb{F}}
\newcommand{\A}{\mathbb{A}}
\newcommand{\PP}{\mathbb{P}}
\newcommand{\E}{\mathbb{E}}
\newcommand{\eps}{\varepsilon}
\newcommand{\dd}{\,\mathrm{d}}
\begin{document}
\section*{UnsolvedMath ID: 701210}
\textbf{Source problem code:} catalogue entry
\subsection{Definitions and mathematical statement}
For a group \(G\), define its derived series by \(G^{(0)}=G\) and \(G^{(i+1)}=[G^{(i)},G^{(i)}]\), where \([H,H]\) is the subgroup generated by the commutators of elements of \(H\). The group \(G\) is soluble of derived length exactly \(n\) if
\[
G^{(n)}=\{1\}\qquad\text{and}\qquad G^{(n-1)}\ne\{1\}.
\]
A total order \(<\) on \(G\) is left invariant if \(x<y\) implies \(gx<gy\) for all \(g,x,y\in G\); it is right invariant if \(x<y\) implies \(xg<yg\). It is invariant on both sides if it has both properties. The supplied record does not specify which invariance convention is intended by ``linearly orderable.''

Fix an intended order-invariance convention, and let \(\mathcal O(G)\) denote the total orders on \(G\) satisfying that convention. Under the usual existential interpretation, a subgroup \(H\le G\) is relatively convex if there exists \(<\in\mathcal O(G)\) such that
\[
\forall h_1,h_2\in H\ \forall g\in G,\qquad
h_1<g<h_2\ \Longrightarrow\ g\in H.
\]
The witnessing order may depend on \(H\). Define
\[
\mathcal R(G)=\{H\le G: H\ne G,\ H\trianglelefteq G,\ H\text{ is relatively convex}\}.
\]
With ``proper'' interpreted literally as \(H\ne G\), the requested existence assertions are
\[
\begin{aligned}
\text{(a)}\quad &\exists G:\quad \mathcal O(G)\ne\varnothing,\quad G^{(3)}=\{1\},\quad G^{(2)}\ne\{1\},\quad |\mathcal R(G)|=1;\\
\text{(b)}\quad &\forall n\in\mathbb Z,\ n>4,\ \exists G:\quad \mathcal O(G)\ne\varnothing,\quad G^{(n)}=\{1\},\quad G^{(n-1)}\ne\{1\},\quad |\mathcal R(G)|=1.
\end{aligned}
\]
Here the group in part (b) may depend on \(n\).

There is an unresolved convention in the uniqueness requirement: the trivial subgroup is proper whenever \(G\ne\{1\}\), and it is convex for every total order. If the source intends ``proper'' to mean nontrivial and proper, then the condition instead becomes
\[
\bigl|\{H\in\mathcal R(G):H\ne\{1\}\}\bigr|=1.
\]
The supplied record does not resolve this distinction or explicitly define relative convexity, so the displayed assertions are conditional formulations rather than a uniquely determined interpretation of the source.

The source additionally states that groups with the intended properties exist for \(n=2\) and \(n=4\). This is retained as a supplied background assertion, without independent verification.
\subsection{Short English statement}
The problem asks whether a linearly orderable soluble group with exactly one proper normal relatively convex subgroup can have derived length \(3\), and whether such a group exists for every integer derived length greater than \(4\). The source reports examples of lengths \(2\) and \(4\). The intended order-invariance convention and whether the unique subgroup must be nontrivial require clarification.
\subsection{Sources}
\begin{itemize}
\item Supplied local UnsolvedMath registry record, permanent ID \(701210\), exact source problem code \texttt{catalogue entry}, titled ``Kourovka Notebook Problem 16.18.'' The record labels the problem open; this is supplied registry status, not an independently verified mathematical assessment.
\item V. V. Bludov, proposer of Problem 16.18 in \emph{The Kourovka Notebook (Unsolved Problems in Group Theory)}, 16th issue (2006), as attributed in the supplied record.
\item The supplied background also identifies the 21st issue edition (2026), \texttt{arXiv:1401.0300v48}. The collection description names E. I. Khukhro and V. D. Mazurov as editors. No source text beyond the supplied record was consulted.
\item Registry provenance: upstream ID \(701210\), upstream revision \texttt{d371ee917ac9647fbb06cd6fc529deec6f3dda82}, and supplied record update timestamp \texttt{2026-10-16T14:00:00Z}. This timestamp is registry metadata, not a literature-verification date.
\end{itemize}
\end{document}
